\subsection{Completeness and round-by-round knowledge soundness of Zinc$+$}
\label{subsec:ucs-pior-security}

We show that the Zinc+ PIOR (\cref{prot:zincplus-ucs-pior}) is complete and round-by-round knowledge sound.

\begin{proof}[Proof of completeness of the Zinc$+$ PIOR (\cref{thm:zincplus-pior-completeness})]
First we show that unless the protocol halts in step 3 or 6 that $(\idx',\inp',\wit') \in \wedge_{\primeindex \in [0,\numprimes]} \ucsalg(\tilde{q}_\primeindex)$. Then we bound the probability of this happening. We begin with $(\idx,\inp,\wit) \in \ucs$, where
\begin{align*}
\idx&=(\witlen,\inplen,\numrows=2^\mu,\primetuple,\bitbound,\degbounds,\constraints = [(Q_{\primeindex\constraintindex}, \langle j_{\primeindex\constraintindex} \rangle)]_{\primeindex \in [0,\numprimes], \constraintindex \in [\numconstraints]}),\enspace
\inp=(\yy)\in(\mainpolyring{\local{\qq}}{\bitbound}{\degreebound_0})^{\inplen} \\
\wit &= (\ff_0,\ff_1,\ldots,\ff_{\numprimes}), \ \ff_0 \in \mainpolyring{\QQ}{\bitbound}{\degreebound_0}^{\witlen},\ \ff_i \in \FF_{q_i}^{<\degreebound_i}[X]^{\witlen}\text{ for } i\in[\numprimes].
\end{align*}
In step~1, the prover sends the oracles for $\tilde{f}_0 \in (\QQ^{<\degreebound_0,\bitbound}[X])[\WWW]$ and $\tilde{f}_{1},\dotsc,\tilde{f}_{\numprimes} \in (\FF^{<\degreebound_i}_{q_i}[X])[\WWW]$, which are the MLEs of the corresponding witness vectors. In step~4, the prover sends
\[
\begin{array}{lll}
  e_{0\constraintindex} &= \sum_{\bb\in\BB^\mu} \phi_{q_0}(Q_{0\constraintindex})(\bb,\phi_{q_0}(\yy),\phi_{q_0}(\ff_0)) \cdot \meq(\bb,\rr_0) & \in \FF_{q_0}[X] \\
  e_{\primeindex\constraintindex} &= \sum_{\bb\in\BB^\mu} \embedding_{\tilde{q}_\primeindex}(Q_{\primeindex\constraintindex})(\bb,\phi_{\tilde{q}_\primeindex}(\yy),\phi_{\tilde{q}_\primeindex}(\ff_0),\embedding_{\tilde{q}_\primeindex}(\ff_\primeindex)) \cdot \meq(\bb,\rr_\primeindex) & \in \FF_{\tilde{q}_\primeindex}[X] \\
\end{array}
\]
for $\primeindex \in [\numprimes], \constraintindex \in [\numconstraints]$. The verifier defines
\[
\begin{array}{lll}
    Q'_{0\constraintindex} = \sum_{\bb\in\BB^\mu} \psi_{q_0,a_0}(Q_{0\constraintindex})(\bb,\inpvar,\ZZZ_1) \cdot \meq(\bb,\rr_0)- e_{0\constraintindex}(a_0) \in \FF_{q_0}[\inpvar,\ZZZ_1] \\
    Q'_{\primeindex\constraintindex} = \sum_{\bb\in\BB^\mu} \psi_{q_i,a_i}(Q_{\primeindex\constraintindex})(\bb,\inpvar,\ZZZ_1,\ZZZ_2) \cdot \meq(\bb,\rr_\primeindex) - e_{\primeindex\constraintindex}(a_i) \in \FF_{\tilde{q}_\primeindex}[\inpvar,\ZZZ_1,\ZZZ_2] \\
\end{array}
\]
Then plugging in $\yy'_i = \psi_{q_i,a_i}(\yy)$, $\ff'_{i0} = \psi_{q_i,a_i}(\ff_{0})$, $\ff'_{i1} = $ we obtain
\[
\begin{array}{rl}
    Q'_{0\constraintindex} =& \sum_{\bb\in\BB^\mu} \left(\psi_{q_0,a_0}(Q_{0\constraintindex})(\bb,\psi_{q_0,a_0}(\yy),\psi_{q_0,a_0}(\ff_{0})) \cdot \meq(\bb,\rr_0) \right)\\
    &- \left(\sum_{\bb\in\BB^\mu} \phi_{q_0}(Q_{0\constraintindex})(\bb,\phi_{q_0}(\yy),\phi_{q_0}(\ff_0)) \cdot \meq(\bb,\rr_0)\right)(a_0) \\
    =& \sum_{\bb\in\BB^\mu} \meq(\bb,\rr_0) \cdot \big[\psi_{q_0,a_0}(Q_{0\constraintindex})(\bb,\psi_{q_0,a_0}(\yy),\psi_{q_0,a_0}(\ff_{0})) \\
    &\phantom{of the opera}\qquad\qquad - \phi_{q_0}(Q_{0\constraintindex})(\bb,\phi_{q_0}(\yy),\phi_{q_0}(\ff_0))(a_0) \big] = 0,
\end{array}
\]
and similarly
\[
\begin{array}{rl}
    Q'_{\primeindex\constraintindex} =& \sum_{\bb\in\BB^\mu} \left(\psi_{q_\primeindex,a_\primeindex}(Q_{\primeindex\constraintindex})(\bb,\psi_{q_\primeindex,a_\primeindex}(\yy),\psi_{q_\primeindex,a_\primeindex}(\ff_{0}), \psi_{q_i,a_i}(\ff_{i})) \cdot \meq(\bb,\rr_\primeindex) \right)\\
    &- \left(\sum_{\bb\in\BB^\mu} \embedding_{\tilde{q}_\primeindex}(Q_{\primeindex\constraintindex})(\bb,\phi_{\tilde{q}_\primeindex}(\yy),\phi_{\tilde{q}_\primeindex}(\ff_0),\embedding_{\tilde{q}_\primeindex}(\ff_\primeindex)) \cdot \meq(\bb,\rr_\primeindex)\right)(a_\primeindex) \\
    =& \sum_{\bb\in\BB^\mu} \meq(\bb,\rr_\primeindex) \cdot \big[\psi_{q_\primeindex,a_\primeindex}(Q_{\primeindex\constraintindex})(\bb,\psi_{q_\primeindex,a_\primeindex}(\yy),\psi_{q_\primeindex,a_\primeindex}(\ff_{\primeindex}), \psi_{q_i,a_i}(\ff_{i})) \\
    &\phantom{of the opera}\qquad\qquad - \embedding_{\tilde{q}_\primeindex}(Q_{\primeindex\constraintindex})(\bb,\phi_{\tilde{q}_\primeindex}(\yy),\phi_{\tilde{q}_\primeindex}(\ff_0),\embedding_{\tilde{q}_\primeindex}(\ff_\primeindex))(a_\primeindex) \big] = 0.
\end{array}
\]
Finally, we conclude our claim of conditional correctness section with the observation that since each $a_i \in \prim(\FF_{\tilde q_i})$ by hypothesis, it holds by~\cref{lem:mle_lifted_projection} that $\incompleteoracle{\tilde{f}'_{\primeindex0}}$ and $\oracle{\tilde{f}'_{\primeindex1}}$ are (complete) polynomial oracles for the MLEs of the corresponding witness vectors.

Now it's left to bound the probability the protocol terminates early. The total number of rational coefficients we apply $\phi_{q_0}$ to is at most $N_{\mathrm{rat}} = \degreebound_0 \cdot (\inplen + \witlen + \numconstraints + \sum_{\constraintindex \in [\numconstraints]} \|\constraintpol_{0\constraintindex} \|_0)$, where $\|\constraintpol \|_0$ denotes the size of the support of $\constraintpol$ as a polynomial in $\rowvar,\inpvar,$ and $\ZZZ_1$. Then using~\cref{lemma:app-prime-division-bound} and the union bound we can conclude the probability $q_0$ divides the denominator of any of these rational coefficients is less than
\[
\frac{N_{\mathrm{rat}} (\bitbound-1)}{|\primeset|(\bitbound_\primeset-1)}.
\]
So given our hypothesis on $|\primeset|(\bitbound_\primeset - 1)$ this happens with probability at most $2^{-\xi}$. Similarly, given our hypothesis on $\extend$ we have that the event $E_i$ that $a_i \not\in \prim(F_{\tilde{q}})$ occurs with probability at most $2^{-\xi}$ (by \cref{lemma:app-primitive-element-sampling}). Taking the union bound over all $E_i$ and also the event $q_0$ divides some denominator gives our completeness bound.
\end{proof}


\begin{proof}[Proof of round-by-round knowledge of the Zinc$+$ PIOR (\cref{thm:zincplus-pior-knowledge})]
We make the simplifying assumption that $(\idx',\inp') \not\in \Relation'$ unless $\incompleteoracle{\tilde{f}'_{00}}$ is complete, which effectively forces $\ff_0 \in \local{q}[X]$. We believe this is justified based on the fact this can be enforced via asking for a single random evaluation of $\incompleteoracle{\tilde{f}'_{00}}$, which will abort with probability greater than $1-\nu/2^{\bitbound_{\primeset}}$ unless complete.

First we define the extractor and knowledge state function, then we upper bound the round-by-round knowledge soundness errors. Throughout we use
\begin{align*}
\idx&=(\witlen=2^\nu,\inplen,\numrows=2^\mu,\primetuple,\bitbound,\degbounds,\constraints = [(Q_{\primeindex\constraintindex}, \langle j_{\primeindex\constraintindex} \rangle)]_{\primeindex \in [0,\numprimes], \constraintindex \in [\numconstraints]}),\enspace
\inp=(\yy)\in(\mainpolyring{\local{\qq}}{\bitbound}{\degreebound_0})^{\inplen},
\end{align*}
and assume $\idx$ is well-formed.

\noindent\textbf{Extractor.} For any transcript $\tr$ where the verifier is about to move, the extractor $\extractor$ queries each polynomial oracle $\oracle{\tilde{f}_i}$ at each point in the $\log m$-dimensional hypercube, and outputs the resulting vectors $(\ff_{0},\ff_1,\dotsc,\ff_{\numprimes})$, where $\ff_{0} \in (\QQ^{<\degreebound_0,\bitbound}[X])^{\witlen}$ and $\ff_{i} \in (\FF_{q_i}^{<\degreebound_i}[X])^{\witlen}\text{ for } i\in[\numprimes]$.

\medskip

\noindent\textbf{Knowledge state function.} We define the knowledge state function's behavior in each round:
\begin{enumerate}[nosep]%\setcounter{enumi}{-1}
    % \item \textbf{Initial transcript:} We set $\kstate(\idx,\inp,\emptyset) = 1$ if and only if the following hold:
    % \begin{itemize}[nosep]
    %     \item Well-formedness: $\kswit = (\ff_0,\ff_1,\ldots,\ff_{\numprimes})$, where $\ff_0 \in (\mainpolyring{\local{\qq}}{\bitbound}{\degreebound_0})^{\witlen}$ and $\ff_i \in (\FF_{q_i}^{<\degreebound_i}[X])^{\witlen}\text{ for } i\in[\numprimes]$
    %     \item Constraint checks: for $\primeindex \in [\numprimes], \constraintindex \in [\numconstraints]$
    %     \[
    %     \begin{array}{rl}
    %     Q_{0 \constraintindex}\left(\rowindex, \yy, \ff_0\right) &\in \langle \generator_{0\constraintindex} \rangle \subseteq \QQ[X] \\
    %     Q_{\primeindex \constraintindex}\left( \rowindex, \phi_{q_\primeindex}(\yy), \phi_{q_\primeindex}(\ff_{0}), \ff_{\primeindex} \right) &\in \langle \generator_{\primeindex\constraintindex} \rangle \subseteq \FF_{q_i}[X].
    %     \end{array}
    %     \]
    % \end{itemize}
    \item \textbf{Ring projection and ideal batching randomness:} At this stage the transcript $\tr$ has the form $\mathrm{tr} = (\oracle{\tilde{f}_0}\|\dotsm\|\oracle{\tilde{f}_{\numprimes}})$. The verifier samples a random prime $q_0\getsr \primeset$ and field elements $\rr_i\getsr \FF_{\tilde q_i}^\mu$ for all $i\in[0,\numprimes]$. We set $\kstate(\idx,\inp,\tr\|q_0\|\rr) = 1$ if and only if running the extractor returns $(\ff_{0},\ff_1,\dotsc,\ff_{\numprimes})$ such that $\ff_0 \in (\ZZ_{(q_0)}[X])^{\witlen}$ and for $\primeindex \in [\numprimes], \constraintindex \in [\numconstraints]$:
    \begin{equation}
    \label{eq:app-ideal-batching}
    \begin{array}{lll}
    \sum_{\bb\in\BB^\mu} \phi_{q_0}\left(Q_{0\constraintindex}(\bb,\yy,\ff_{0})\right) \cdot \meq(\bb,\rr_0) &\in \langle \phi_{q_0}(\generator_{0\constraintindex}) \rangle &\subseteq \FF_{q_0}[X], \\
    \sum_{\bb\in\BB^\mu} \embedding_{\tilde{q}_\primeindex}(Q_{\primeindex\constraintindex})(\bb,\phi_{\tilde{q}_\primeindex}(\yy),\phi_{\tilde{q}_i}(\ff_{0}),\phi_{\tilde{q}_i}(\ff_{i})) \cdot \meq(\bb,\rr_\primeindex) &\in \langle \embedding_{\tilde{q}_i}(\generator_{\primeindex\constraintindex}) \rangle &\subseteq \FF_{\tilde{q}_i}[X].
    \end{array}
    \end{equation}
    Recall that since the index is well-formed this also implies $\ff_0 \in (\ZZ_{(q_1,\dotsc,q_{\numprimes})}[X])^{\witlen}$.

    \item \textbf{Field projection randomness:} At this stage the transcript $\tr$ has the form $\mathrm{tr} = (\oracle{\tilde{f}_0}\|\dotsm\|\oracle{\tilde{f}_{\numprimes}},q_0\|\rr,[e_{\primeindex\constraintindex}]_{\primeindex \in [0,\numprimes],\constraintindex \in [\numconstraints]})$. The verifier samples random field elements $a_i \getsr \FF_{\tilde q_i}$, for $i\in[0,\numprimes]$. We set $\kstate(\idx,\inp,\tr\|q_0\|\rr) = 1$ if and only if running the extractor returns $(\ff_{0},\ff_1,\dotsc,\ff_{\numprimes})$ such that $\ff_0 \in (\ZZ_{(q_0)}[X])^{\witlen}$ and for $\primeindex \in [\numprimes], \constraintindex \in [\numconstraints]$:
    \begin{equation}
        \label{eq:app-kstate-projection}
    \begin{array}{ll}
    \left(\sum_{\bb\in\BB^\mu} \phi_{q_0}(Q_{0\constraintindex})(\bb,\phi_{q_0}(\yy),\phi_{q_0}(\ff_{0})) \cdot \meq(\bb,\rr_0) - e_{0\constraintindex} \right)(a_0) = 0 & \textnormal{over } \FF_{q_0} \\
    \left(\sum_{\bb\in\BB^\mu} \embedding_{\tilde{q}_\primeindex}(Q_{\primeindex\constraintindex})(\bb,\phi_{\tilde{q}_\primeindex}(\yy),\phi_{\tilde{q}_i}(\ff_{0}),\phi_{\tilde{q}_\primeindex}(\ff_{\primeindex})) \cdot \meq(\bb,\rr_\primeindex) - e_{\primeindex\constraintindex} \right)(a_\primeindex) = 0 & \textnormal{over } \FF_{\tilde{q}_\primeindex}
    \end{array}
    \end{equation}
    
    \item \textbf{Full transcript:} It remains to show for a full transcript $\tr=(\oracle{\tilde{f}_0}\|\dotsm\|\oracle{\tilde{f}_{\numprimes}},q_0\|\rr,\mathbf{e},\mathbf{a})$ that $\kstate(\idx,\inp,\tr)=1$ if and only if $\verifier\bigl(\idx,\inp,\tr\bigr)$ outputs $(\idx',\inp')$ such that there exists $\wit$' with $(\idx',\inp',\wit')\in \Relation'$.
    
    First assume $\verifier\bigl(\idx,\inp,\tr\bigr)$ outputs $(\idx',\inp')$ such that there exists $\wit$' with $(\idx',\inp',\wit')\in \Relation'$. Since $(\idx',\inp') \not\in \Language(\Relation')$ unless $\incompleteoracle{\tilde{f}'_{00}}$ is complete by assumption, we have that $\ff_0 \in \local{q_0}$. The $\wit$' is by definition of UCS and by~\cref{lem:mle_lifted_projection} comprised of the vectors obtained by applying $\psi$ to the $\ff_i$ as in an honest run of the protocol. Then ~\cref{eq:app-kstate-projection} is exactly checking if that unique $\wit$' satisfies $(\idx',\inp',\wit')\in \Relation'$, outputting $\kstate(\idx,\inp,\tr)=1$ if so.
    
    Next assume $\kstate(\idx,\inp,\tr)=1$. In this case it is straightforward to see we can explicitly produce a $\wit$' such that $(\idx',\inp',\wit')\in \Relation'$ using the extracted vectors.
\end{enumerate}

\medskip
\noindent\textbf{RBR knowledge soundness.} We upper bound the round-by-round knowledge soundness errors based on the knowledge state function and extractor described above.
\begin{enumerate}[nosep]
    \item \textbf{Ring projection and ideal batching randomness:} We show that
    \[
    \Pr_{q_0,\rr}\left[
    \begin{array}{l}
    \kstate(\idx,\inp,(\oracle{f_i})_{i =0}^{\numprimes}\Vert(q_0,\rr))=1 \\
    \wedge\ (\idx,\inp,\extractor(\idx,\inp,(\oracle{f_i})_{i =0}^{\numprimes})) \not\in \Relation
    \end{array}
    \right]
    \ \le\ 2^{-\lambda}.
    \]
    By~\cref{lemma:app-phi-collision} we have that
    \[
    \Pr_{q_0 \gets \primeset}\left[\phi_{q_0}\left(\constraintpol_{0\constraintindex}(\rowindex,\yy,\ff_0)\right) \in \langle \phi_{q_0}(\generator_{0\constraintindex}) \rangle \ \wedge\  \constraintpol_{0\constraintindex}(\rowindex,\yy,\ff_0) \not\in \langle \generator_{0\constraintindex} \rangle\right] < \frac{8\degreebound_{0}^{2}\bitbound}{|\primeset| (\bitbound_\primeset-1)} < 2^{-\lambda-1}.
    \]
    We note applying the canonical embedding cannot result in an unsatisfied constraint over $\FF_{q_i}$ being satisfied in $\FF_{\tilde{q}_i}$.
    
    Else suppose for some $(k, \primeindex, \constraintindex)$ a constraint doesn't hold over $\FF_{q_0}$ or $\FF_{\tilde{q}_i}$. In this case~\cref{lem:ideal-batching} tells us that~\cref{eq:app-ideal-batching} will hold with probability at most $\mu/|\FF_{\tilde{q}_i}| \leq 2^{-\lambda-1}$. A union bound over these two cases gives $\knowledgeerror_1 = 2^{-\lambda}$.
    
    We now turn to the second source of randomness, namely the field-projection points~$\evalpoints$, and bound the probability that the knowledge state transitions from~$0$ to~$1$ in this round.
    
    \item \textbf{Field projection randomness:} We show that for all transcripts of the form $\tr = (\oracle{f_i})_{i =0}^{\numprimes}\|(q_0,\rr)\|[e_{\primeindex\constraintindex}]_{\primeindex\constraintindex}$ that
    \[
    \Pr_{\evalpoints}\left[
    \begin{array}{l}
    \kstate(\idx,\inp,\tr) =0 \\
    \wedge\ \kstate(\idx,\inp,\tr\|\evalpoints)=1 \\
    \wedge\ (\idx,\inp,\extractor(\idx,\inp,\tr)) \not\in \Relation
    \end{array}
    \right]
    \ \le\ 2^{-\lambda}.
    \]
    To do this we must consider the case that for some $\primeindex \in [0,\numprimes], \constraintindex \in [\numconstraints]$ \cref{lem:ideal-batching} doesn't hold, but when projected to $\FF_{\tilde{q}_i}$ via $\psi_{a_i}$ it does. This happens with probability at most $2 m \ell d_i \mathsf{deg}(\constraintpol_{it}) /|\prim(\FF_{\tilde{q}_i})| \leq 2^{-\lambda} = \knowledgeerror_2$.
\end{enumerate}
\end{proof}

\parhead{Supplementary lemmas}
Here we include some lemmas we draw on in our proofs of completeness and RBR knowledge.

\begin{lemma}
    \label{lemma:app-prime-division-bound}
    Let $\bitbound > \bitbound_\primeset$. For any $a \in \ZZ^{<\bitbound} \setminus \{0\}$ of bit length less than $\bitbound$, the number of primes of bit length at least $B_\primeset$ dividing it is less than $(\bitbound-1)/(\bitbound_\primeset-1)$.
\end{lemma}
\begin{proof}
Say $t$ primes $p_1,\dotsc,p_t$ divide $a$. By assumption $|a| < 2^{\bitbound-1}$ and that $p_i \geq 2^{\bitbound_\primeset-1}$ for all $i \in [t]$. Then we have $2^{t(\bitbound_\primeset-1)} \leq |a| < 2^{\bitbound-1}$, which implies $t < (\bitbound-1)/(\bitbound_\primeset-1)$.
\end{proof}

\begin{lemma}
\label{lemma:app-phi-collision}
Let $f,\generator \in \QQ^{<\degreebound,<\bitbound}[X]$, $\degreebound \ge 1$, $\bitbound > \bitbound_\primeset \ge 2$.
Then
\[
\Pr_{p \getsr \primeset}\left[\phi_p(f) \in \langle \phi_p(\generator) \rangle \ \wedge\ f \not\in \langle \generator \rangle \pST f,\generator \in \ZZ_{(p)}\right]
< \frac{8 d^2 B}{|\primeset|(\bitbound_\primeset-1)}.
\]
\end{lemma}
\begin{proof}
Let $(q,r) \in \QQ[X]\times \QQ[X]$ be the unique quotient--remainder pair given by Euclidean division in $\QQ[X]$:
\[
f = \generator q + r,
\qquad
\deg(r) < \deg(\generator).
\]
Then $f \in \langle \generator\rangle$ if and only if $r=0$.

Write $r$ in lowest terms as $r=n_r/d_r$, where $n_r\in \ZZ[X]$ and $d_r\in \ZZ^{+}$ satisfy $\gcd(d_r,c)=1$, where $c = \mathsf{cont}(n_r)$ is the content of $n_r$. Similarly write $a$, the leading coefficient of $\generator$, in lowest terms $a = n_a/d_a$. Define $N := c \cdot n_a$. We show that for every prime $p$,
\[
\left(\phi_p(f) \in \langle \phi_p(\generator) \rangle \ \wedge\ f \not\in \langle \generator \rangle \ \wedge\ f,\generator \in \ZZ_{(p)}\right)
\Longrightarrow p \mid N.
\]
Fix $p$ and assume $\phi_p(f) \in \langle \phi_p(\generator) \rangle$, $f \not\in \langle \generator \rangle$, and $f,\generator \in \ZZ_{(p)}$.
Define $N := c \cdot n_a$. We show $p \mid N$. First assume $p \nmid n_a$. Then the leading coefficient of $\generator$ is a unit in $\ZZ_{(p)}$, so Euclidean division of $f$ by $\generator$ can be performed in $\ZZ_{(p)}[X]$. Let $(q_p,r_p)\in \ZZ_{(p)}[X]\times \ZZ_{(p)}[X]$ be the resulting quotient--remainder pair, so
\[
f = \generator q_p + r_p,
\qquad
\deg(r_p) < \deg(\generator).
\]
By uniqueness of Euclidean division in $\QQ[X]$, we have $r_p=r$, and in particular $r\in \ZZ_{(p)}[X]$. Applying $\phi_p$ gives
\[
\phi_p(f) = \phi_p(\generator)\cdot \phi_p(q_p) + \phi_p(r).
\]
Since $p \nmid n_a$, we have $\deg(\phi_p(\generator))=\deg(\generator)$, hence $\deg(\phi_p(r))<\deg(\phi_p(\generator))$.  Together with $\phi_p(f)\in \langle \phi_p(\generator)\rangle$, this implies $\phi_p(r)=0$. Because $r\in \ZZ_{(p)}[X]$, we have $\phi_p(r)=0$ if and only if $p \mid \mathsf{cont}(n_r)=c$. Hence $p \mid c$, and therefore $p \mid N$.

Next assume $p \nmid c$. If also $p \nmid n_a$, then the same argument shows $r\in \ZZ_{(p)}[X]$ and $\phi_p(r)=0 \Longleftrightarrow p \mid c$, so $p \nmid c$ implies $\phi_p(r)\neq 0$ while $\deg(\phi_p(r))<\deg(\phi_p(\generator))$. A nonzero polynomial of degree strictly less than $\deg(\phi_p(\generator))$ cannot lie in $\langle \phi_p(\generator)\rangle$, contradicting $\phi_p(f)\in \langle \phi_p(\generator)\rangle$.  Therefore $p \mid n_a$, and hence $p \mid N$.

Let $\primeset_\textnormal{bad} = \{p \in \primeset : p \mid N\}$. The above implication yields
\[
\Pr_{p \getsr \primeset}\left[\phi_p(f) \in \langle \phi_p(\generator) \rangle \ \wedge\ f \not\in \langle \generator \rangle \pST f,\generator \in \ZZ_{(p)}\right]
\le \frac{|\primeset_\textnormal{bad}|}{|\primeset|}.
\]
Let $\bitbound_N = \lceil \log(|N|+1) \rceil$. If $\bitbound_N \le \bitbound_\primeset$ then $|\primeset_\textnormal{bad}|=0$.
Otherwise, by~\cref{lemma:app-prime-division-bound},
\[
|\primeset_\textnormal{bad}| < \frac{\bitbound_N-1}{\bitbound_\primeset-1},
\]
so
\begin{equation}
\label{eq:app-abstract-phi-collision-bound}
\Pr_{p \getsr \primeset}[\cdots]
< \frac{\bitbound_N-1}{|\primeset|(\bitbound_\primeset-1)}.
\end{equation}

We now bound $\bitbound_N$. Since the coefficients of $f$ and $\generator$ have bitlength $<B$, we may write each coefficient in lowest
terms with numerator and denominator of magnitude $<2^B$. Let $\Delta\in \ZZ^{+}$ be the product of all coefficient denominators appearing
in $f$ and $\generator$. Then $\Delta < 2^{2dB}$, and the polynomials
\[
F := \Delta f \in \ZZ[X], \qquad J := \Delta \generator \in \ZZ[X]
\]
satisfy $\|F\|_\infty,\|J\|_\infty < 2^B \Delta < 2^{(2d+1)B}$.

Let $m := \deg(F)$ and $n := \deg(J)$, so $m,n < d$. Let $b$ be the leading coefficient of $J$. By pseudo-division, there exist
$Q',R' \in \ZZ[X]$ with $\deg(R')<n$ such that
\begin{equation}\label{eq:app-pseudodiv_finish}
b^{m-n+1} F = J Q' + R'.
\end{equation}
Dividing \eqref{eq:app-pseudodiv_finish} by $\Delta$ and comparing with $f=\generator q + r$ shows that
\[
r = \frac{R'}{b^{m-n+1}\Delta}.
\]
In particular, there exists $\alpha\in\ZZ^{+}$ such that $R'=\alpha n_r \in \ZZ[X]$ and $b^{m-n+1}\Delta=\alpha d_r \in \ZZ\setminus\{0\}$,
hence
\[
c = \mathsf{cont}(n_r) \le \|n_r\|_\infty \le \|R'\|_\infty.
\]
Each coefficient of $R'$ can be expressed as a determinant of a matrix of dimension at most $m+n < 2d$ whose entries are coefficients of $F$
and $J$. By Hadamard's inequality,
\[
\|R'\|_\infty \le (2d)^d \cdot \bigl(\max(\|F\|_\infty,\|J\|_\infty)\bigr)^{2d}
< (2d)^d \cdot \left(2^{(2d+1)B}\right)^{2d}
= (2d)^d \cdot 2^{(4d^2+2d)B}.
\]
Therefore $c < (2d)^d \cdot 2^{(4d^2+2d)B}$. Since $|n_a|<2^B$, we have
\[
|N| < (2d)^d \cdot 2^{(4d^2+2d+1)B}.
\]
It follows that
\[
\bitbound_N = \lceil \log(|N|+1) \rceil \le (4d^2+2d+1)B + d\log(2d) + 1 < 8 d^2 B.
\]
Substituting this bound into~\cref{eq:app-abstract-phi-collision-bound} completes the proof.
\end{proof}

\begin{lemma}
\label{lemma:app-primitive-element-sampling}
Let $q$ be a prime power, $k\ge 2$ an integer. Let $\extend(q) = q^k$ pick $k$ such that $q^k \geq 2^{2\xi} \cdot (8(\xi+1))^2$. Define $\prim(\FF_{q^k}) = \{\alpha\in \FF_{q^k} \ \mid \ \FF_q(\alpha)=\FF_{q^k}\}$. Then
\[
\frac{|\prim(\FF_{q^k})|}{q^k} \ge 1-2^{-\xi}.
\]
\end{lemma}
\begin{proof}
An element $\alpha\in \FF_{q^k}$ fails to generate $\FF_{q^k}$ over $\FF_q$ if and only if
$\alpha$ lies in a proper intermediate subfield $\FF_{q^d}$ for some divisor $d\mid k$ with $d<k$.
Thus
\[
\FF_{q^k}\setminus \prim(\FF_{q^k}) \subseteq \bigcup_{\substack{d\mid k\\ d<k}} \FF_{q^d}.
\]
Taking cardinalities and applying the union bound yields
\[
|\FF_{q^k}\setminus \prim(\FF_{q^k})|
\le \sum_{\substack{d\mid k\\ d<k}} |\FF_{q^d}|
= \sum_{\substack{d\mid k\\ d<k}} q^d.
\]
Dividing by $q^k$ gives
\[
1-\frac{|\prim(\FF_{q^k})|}{q^k}
\le \sum_{\substack{d\mid k\\ d<k}} q^{d-k}.
\]
Let $\ell$ denote the smallest prime divisor of $k$. Then every proper divisor $d$ of $k$ satisfies
$d\le k/\ell$, so $q^{d-k}\le q^{k/\ell-k}=q^{-(k-k/\ell)}$. Since $k$ has at most $k-1$ proper divisors,
\[
1-\frac{|\prim(\FF_{q^k})|}{q^k}
\le (k-1)\,q^{-(k-k/\ell)}.
\]
Using $\ell\ge 2$, we have $k-k/\ell \ge k/2$, hence
\[
(k-1)\,q^{-(k-k/\ell)} \le (k-1)\,q^{-k/2}.
\]
Write $N := q^k$. Then $q^{-k/2}=N^{-1/2}$, so the previous equation becomes
\[
1-\frac{|\prim(\FF_{q^k})|}{q^k}
\le (k-1)\,N^{-1/2}.
\]
Moreover, since $q\ge 2$ and $k\ge 2$, we have $N=q^k\ge 2^k$, and thus $k-1\le \log_2 N$.
Therefore
\[
1-\frac{|\prim(\FF_{q^k})|}{q^k}
\le \frac{\log_2 N}{\sqrt{N}}.
\]
By the hypothesis $N=q^k \ge 2^{2\xi}\cdot (8(\xi+1))^2$, we have
$\sqrt{N}\ge 2^\xi\cdot 8(\xi+1)$. Also,
\[
\log_2 N
\le \log_2\!\Big(2^{2\xi}\cdot (8(\xi+1))^2\Big)
= 2\xi + 2\log_2(8(\xi+1))
\le 8(\xi+1),
\]
where the last inequality uses $\xi\ge 1$ and $\log_2(8(\xi+1))\le 3+(\xi+1)$.
Combining the above bounds gives
\[
\frac{\log_2 N}{\sqrt{N}}
\le \frac{8(\xi+1)}{2^\xi\cdot 8(\xi+1)}
= 2^{-\xi}.
\]
Hence $1-\frac{|\prim(\FF_{q^k})|}{q^k}\le 2^{-\xi}$, which is equivalent to the claimed bound.
\end{proof}